% GENERATED FILE. Edit canonical lessons, then run npm run generate:books.
% canonical-source-hash: fnv1a64:24d6742864679dd8
% canonical-lessons: ML-C307-valarttuka, ML-C307-valaruka, ML-C307-abhinayikkuka, ML-C307-pottikkuka, ML-C307-pukalttuka

\chapter{To raise, To grow, To act, To burst, To praise}
\label{ch:ml-to-raise-to-grow-to-act-to-burst-to-praise}
\hlchaptermodality{307}

\begin{chapteropening}
\textbf{By the end of this chapter:} \emph{I can say to raise, to grow, to act, to burst and to praise.}

Complete the last lesson of chapter 307: I can say to raise, to grow, to act, to burst and to praise.
\end{chapteropening}

\section[\texorpdfstring{\ml{വളർത്തുക} (vaḷarttuka)}{വളർത്തുക (vaḷarttuka)}]{\ml{വളർത്തുക} (vaḷarttuka) — to raise, to rear (children, pets)}

\label{lesson:ML-C307-valarttuka}



\begin{quote}
Before the new one: say the Malayalam for to postpone, then the Malayalam for to feel.
\end{quote}

\subsection*{You'll want to know: \ml{വളർത്തുക}}
\textbf{\ml{വളർത്തുക}} — \emph{vaḷarttuka} — \textquotedblleft{}to raise, to rear (children, pets)\textquotedblright{}.

\begin{grammarlens}[title={What you've built}]
One more word for this chapter.
\end{grammarlens}

\subsection*{Your turn}
\begin{itemize}
\raggedright
  \item \emph{Say it:} \emph{vaḷarttuka}
  \item \emph{Say it:} \emph{vaḷarttuka}, once more
  \item \emph{Say it:} say \emph{vaḷarttuka}
  \item \emph{Recall:} say the Malayalam for to postpone, then the Malayalam for to feel, then say \emph{vaḷarttuka} again
\end{itemize}

\begin{tcolorbox}[breakable,colback=teal!4,colframe=teal!35!black,title={Before you move on}]
What does \ml{വളർത്തുക} mean? (\textquotedblleft{}To raise, to rear (children, pets)\textquotedblright{}.) Say it once more.
\end{tcolorbox}
\section[\texorpdfstring{\ml{വളരുക} (vaḷaruka)}{വളരുക (vaḷaruka)}]{\ml{വളരുക} (vaḷaruka) — to grow}

\label{lesson:ML-C307-valaruka}



\begin{quote}
Before the new one: say the Malayalam for to feel, then the Malayalam for to raise.
\end{quote}

\subsection*{You'll want to know: \ml{വളരുക}}
\textbf{\ml{വളരുക}} — \emph{vaḷaruka} — \textquotedblleft{}to grow\textquotedblright{}.

\begin{grammarlens}[title={What you've built}]
One more word for this chapter.
\end{grammarlens}

\subsection*{Your turn}
\begin{itemize}
\raggedright
  \item \emph{Say it:} \emph{vaḷaruka}
  \item \emph{Say it:} \emph{vaḷaruka}, once more
  \item \emph{Say it:} say \emph{vaḷaruka}
  \item \emph{Recall:} say the Malayalam for to feel, then the Malayalam for to raise, then say \emph{vaḷaruka} again
\end{itemize}

\begin{tcolorbox}[breakable,colback=teal!4,colframe=teal!35!black,title={Before you move on}]
What does \ml{വളരുക} mean? (\textquotedblleft{}To grow\textquotedblright{}.) Say it once more.
\end{tcolorbox}
\section[\texorpdfstring{\ml{അഭിനയിക്കുക} (abhinayikkuka)}{അഭിനയിക്കുക (abhinayikkuka)}]{\ml{അഭിനയിക്കുക} (abhinayikkuka) — to act (in a play or film)}

\label{lesson:ML-C307-abhinayikkuka}



\begin{quote}
Before the new one: say the Malayalam for to raise, then the Malayalam for to grow.
\end{quote}

\subsection*{You'll want to know: \ml{അഭിനയിക്കുക}}
\textbf{\ml{അഭിനയിക്കുക}} — \emph{abhinayikkuka} — \textquotedblleft{}to act (in a play or film)\textquotedblright{}.

\begin{grammarlens}[title={What you've built}]
One more word for this chapter.
\end{grammarlens}

\subsection*{Your turn}
\begin{itemize}
\raggedright
  \item \emph{Say it:} \emph{abhinayikkuka}
  \item \emph{Say it:} \emph{abhinayikkuka}, once more
  \item \emph{Say it:} say \emph{abhinayikkuka}
  \item \emph{Recall:} say the Malayalam for to raise, then the Malayalam for to grow, then say \emph{abhinayikkuka} again
\end{itemize}

\begin{tcolorbox}[breakable,colback=teal!4,colframe=teal!35!black,title={Before you move on}]
What does \ml{അഭിനയിക്കുക} mean? (\textquotedblleft{}To act (in a play or film)\textquotedblright{}.) Say it once more.
\end{tcolorbox}
\section[\texorpdfstring{\ml{പൊട്ടിക്കുക} (poṭṭikkuka)}{പൊട്ടിക്കുക (poṭṭikkuka)}]{\ml{പൊട്ടിക്കുക} (poṭṭikkuka) — to burst, to set off (firecrackers)}

\label{lesson:ML-C307-pottikkuka}



\begin{quote}
Before the new one: say the Malayalam for to grow, then the Malayalam for to act.
\end{quote}

\subsection*{You'll want to know: \ml{പൊട്ടിക്കുക}}
\textbf{\ml{പൊട്ടിക്കുക}} — \emph{poṭṭikkuka} — \textquotedblleft{}to burst, to set off (firecrackers)\textquotedblright{}.

\begin{grammarlens}[title={What you've built}]
One more word for this chapter.
\end{grammarlens}

\subsection*{Your turn}
\begin{itemize}
\raggedright
  \item \emph{Say it:} \emph{poṭṭikkuka}
  \item \emph{Say it:} \emph{poṭṭikkuka}, once more
  \item \emph{Say it:} say \emph{poṭṭikkuka}
  \item \emph{Recall:} say the Malayalam for to grow, then the Malayalam for to act, then say \emph{poṭṭikkuka} again
\end{itemize}

\begin{tcolorbox}[breakable,colback=teal!4,colframe=teal!35!black,title={Before you move on}]
What does \ml{പൊട്ടിക്കുക} mean? (\textquotedblleft{}To burst, to set off (firecrackers)\textquotedblright{}.) Say it once more.
\end{tcolorbox}
\section[\texorpdfstring{\ml{പുകഴ്ത്തുക} (pukaḻttuka)}{പുകഴ്ത്തുക (pukaḻttuka)}]{\ml{പുകഴ്ത്തുക} (pukaḻttuka) — to praise}

\label{lesson:ML-C307-pukalttuka}



\begin{quote}
Before the new one: say the Malayalam for to act, then the Malayalam for to burst.
\end{quote}

\subsection*{You'll want to know: \ml{പുകഴ്ത്തുക}}
\textbf{\ml{പുകഴ്ത്തുക}} — \emph{pukaḻttuka} — \textquotedblleft{}to praise\textquotedblright{}.

\begin{grammarlens}[title={What you've built}]
One more word for this chapter.
\end{grammarlens}

\subsection*{Your turn}
\begin{itemize}
\raggedright
  \item \emph{Say it:} \emph{pukaḻttuka}
  \item \emph{Say it:} \emph{pukaḻttuka}, once more
  \item \emph{Say it:} say \emph{pukaḻttuka}
  \item \emph{Recall:} say the Malayalam for to act, then the Malayalam for to burst, then say \emph{pukaḻttuka} again
\end{itemize}

\begin{tcolorbox}[breakable,colback=teal!4,colframe=teal!35!black,title={Before you move on}]
What does \ml{പുകഴ്ത്തുക} mean? (\textquotedblleft{}To praise\textquotedblright{}.) Say it once more.
\end{tcolorbox}
